\section{Introducción: ¿qué es un LLM Agent?}

Shunyu Yao (姚顺雨) es el creador del marco ReAct y actualmente trabaja en OpenAI. Esta clase ofrece un recuento sistemático de la definición de los LLM Agent, la trayectoria histórica de su desarrollo y las direcciones futuras.

\subsection{La definición de Agent}

\begin{importantbox}{Las tres capas de la definición de Agent}
\begin{enumerate}
    \item \textbf{Text Agent}: un sistema inteligente que interactúa con el entorno, cuyas acciones (action) y observaciones (observation) se expresan en forma de lenguaje (text). El Text Agent no necesariamente usa un LLM; desde los inicios mismos de la IA ya existían Text Agent basados en reglas.
    \item \textbf{LM Agent}: un Text Agent que usa un modelo de lenguaje para generar sus acciones.
    \item \textbf{Reasoning Agent}: un Agent que usa un modelo de lenguaje para \textbf{razonar antes de actuar}, lo cual constituye un salto cualitativo.
\end{enumerate}
\end{importantbox}

Definir Agent requiere responder dos subpreguntas: ¿qué es la ``inteligencia'' (intelligence)? ¿Qué es el ``entorno'' (environment)?

\begin{itemize}
    \item El entorno puede ser físico (robots, conducción autónoma) o digital (videojuegos, páginas web, editores de código).
    \item La definición de ``inteligencia'' evoluciona con el tiempo: hace 60 años, un simple chatbot ya se consideraba inteligente, mientras que hoy ChatGPT ya no sorprende a nadie.
\end{itemize}

\subsection{Resumen de la sección}

La innovación central del LLM Agent no consiste en hacer que el modelo de lenguaje ejecute acciones, sino en hacer que \textbf{razone antes de actuar}. Lo que distingue al Reasoning Agent de todos los paradigmas de Agent anteriores es que el razonamiento (thinking) es una acción interna con un espacio de lenguaje ilimitado.
